\section{图形处理器的起源、统一计算设备架构与应用扩张}

\subsection{图形工作负载天然具有大规模数据并行性}

图形处理器最初为图形渲染而设计，因此名称中的“图形”并不是历史偶然。屏幕包含大量像素，许多像素可以在相互独立或弱依赖的条件下执行相近的颜色、几何与纹理计算。把相同计算作用于大量不同数据元素，正符合高吞吐并行硬件的设计取向。

2007 年以前，图形处理器主要通过图形应用程序编程接口使用，例如\term{开放图形库}{OpenGL}与\term{直接三维图形接口}{Direct3D}。研究者若想让图形处理器完成非图形数值计算，必须把问题重新表述为图形流水线能够接受的形式，即把一般计算改写成“为某个像素着色的函数”，再把渲染输出解释为所需数值结果。

这种方法证明了图形处理器具有通用数值计算潜力，但编程过程高度间接，算法表达受图形接口约束。它促使硬件与软件接口共同演化。

\begin{figure}[H]
  \centering
  \includegraphics[width=0.91\textwidth]{gpu_origins.png}
  \caption{图形处理器从图形接口走向通用计算的历史转折}
\end{figure}

\subsection{2007 年的统一计算设备架构转折}

2007 年，英伟达发布\term{统一计算设备架构}{Compute Unified Device Architecture, CUDA}。它为图形处理器提供面向通用计算的编程接口，使程序员可以直接表达线程、内核函数、设备内存与数据传输，而不必把数值问题伪装成图形渲染。

这一转折通常称为\term{通用图形处理器计算}{General-Purpose Computing on Graphics Processing Units, GPGPU}。它包含两方面变化：
\begin{itemize}
  \item 软件层面提供适合 C/C++ 程序员使用的语言扩展、运行时接口与工具链；
  \item 硬件层面增加或调整支持通用加载、存储、控制与数值计算的结构，使图形处理器不再只适合固定图形流水线。
\end{itemize}

统一计算设备架构 C/C++ 程序通常同时包含中央处理器代码和图形处理器代码：中央处理器端创建数据、分配设备资源并发起计算；图形处理器端执行由大量线程组成的\term{内核函数}{kernel}。这种主机端与设备端的协作构成统一计算设备架构程序的高层执行模型。

\subsection{为什么沿用图形处理器，而不是从零设计另一种吞吐处理器}

先进芯片的设计、验证、制造与软件生态建设成本极高，需要足够大的销量分摊固定成本。完全新建一种大规模并行处理器，即使技术上可行，也必须同时建立市场、编译器、工具、库和开发者生态。

图形处理器在并行计算进入主流之前已经拥有庞大的游戏与图形市场。这个既有装机基础带来三项先发优势：
\begin{enumerate}
  \item 规模化销量能够分摊芯片研发与制造成本；
  \item 图形工作负载已经推动硬件形成大规模并行执行能力；
  \item 既有产品迭代、驱动与开发工具可以逐步扩展到科学计算和人工智能，而不必从零建立完整生态。
\end{enumerate}

因此，图形处理器的成功来自技术结构与市场结构的共同作用。它既具备适合吞吐计算的硬件基础，也拥有其他专用加速器往往缺乏的持续商业规模。

\subsection{数据中心与人工智能推动通用图形处理器计算扩张}

英伟达 2019--2024 年按市场部门划分的收入曲线反映了通用图形处理器计算的扩张。图中数据中心收入在后期快速上升，这一变化与人工智能驱动的训练和推理需求密切相关。

\begin{figure}[H]
  \centering
  \includegraphics[width=0.83\textwidth]{gpgpu_revenue_growth.png}
  \caption{数据中心收入增长与通用图形处理器计算的扩张}
\end{figure}

该图反映了应用重心的变化：图形处理器仍服务于游戏与可视化，但数据中心、高性能计算和人工智能已成为其重要使用场景。模型训练需要执行大规模矩阵与张量运算；模型推理则把类似计算转化为持续服务需求。两者都可产生规模巨大且结构相对规则的并行工作。

\subsection{人工智能计算需求的数量级增长}

若干机器学习系统所需的浮点运算数量表明，从早期模型到大型语言与多模态模型，计算量跨越了许多数量级，最高约为 $10^{24}$ 次浮点运算。

\begin{figure}[H]
  \centering
  \includegraphics[width=0.92\textwidth]{ai_compute_needs.png}
  \caption{人工智能系统计算需求增长趋势}
\end{figure}

对数坐标上的纵向小幅移动代表乘法级增长。例如从 $10^{16}$ 增长到 $10^{20}$ 不是增加四个单位，而是增加 $10^4$ 倍。大模型计算需求的这种数量级变化，使单纯依靠单核性能提升不可行，也使高吞吐加速器、多设备并行与大规模系统调度成为必要条件。

\subsection{超级计算机中的加速器普及}

全球超级计算机五百强榜单依据大规模线性代数基准对系统进行排序，通常每半年更新一次，分别对应年中和年末的主要超级计算会议。2026 年 6 月榜单的前五名中，第一名使用自定义多核处理器；第二至第五名分别明确使用 AMD Instinct MI300A、AMD Instinct MI250X、英特尔数据中心图形处理器 Max 和英伟达 GH200 超级芯片。这个样本说明，顶级系统通常依靠大量加速器获得主要浮点计算能力。

\begin{figure}[H]
  \centering
  \includegraphics[width=0.98\textwidth]{top500_gpu_prevalence.png}
  \caption{2026 年 6 月全球超级计算机五百强前五名}
\end{figure}

中央处理器仍负责操作系统、作业管理、控制流和部分串行计算；加速器负责高密度并行数值工作。要让几十年前为串行处理器编写的科学代码利用这些系统，必须重构算法、数据布局和并行分解。硬件峰值只是潜在能力，应用性能来自软件对硬件的有效映射。

\subsection{并行编程框架之间的关系}

统一计算设备架构是并行编程体系中的一种平台；其他常见框架分别提供跨设备编程、指令式卸载、分布式通信或共享内存并行等抽象。

\begin{table}[H]
\centering
\caption{常见并行编程框架}
\begin{tabularx}{\textwidth}{L{0.22\textwidth}L{0.30\textwidth}Y}
\toprule
\textbf{框架} & \textbf{主要目标} & \textbf{与统一计算设备架构的关系} \\
\midrule
统一计算设备架构 & 英伟达图形处理器的通用并行编程平台 & 本课程的主要编程环境。 \\
\term{开放计算语言}{OpenCL} & 面向多类处理器与加速器的开放并行编程接口 & 提供另一类图形处理器并行编程框架。 \\
\term{开放加速器编程指令}{OpenACC} & 通过编译指令把循环与数据区域卸载到加速器 & 可用于较高层次的加速器编程，本课程不以它为主线。 \\
\term{消息传递接口}{Message Passing Interface, MPI} & 分布式内存系统中进程之间的显式通信 & 可与统一计算设备架构组合，在多个计算节点和多张图形处理器之间分配工作。 \\
\term{开放式多处理}{Open Multi-Processing, OpenMP} & 共享内存系统中的线程级并行与编译指令 & 可与统一计算设备架构组合，协调单节点内的中央处理器线程及多设备工作。 \\
\bottomrule
\end{tabularx}
\end{table}

单张图形处理器只解决一个设备内部的并行执行。若应用要扩展到多个节点，就需要额外机制完成进程间数据交换、任务划分和同步；消息传递接口承担这类分布式协调。若应用在单个共享内存节点内使用多个中央处理器线程或多个加速器，开放式多处理可以参与上层协调。统一计算设备架构、消息传递接口和开放式多处理因此可以处于同一应用的不同层次，而不是相互排斥的替代方案。

\begin{keybox}
图形处理器成为通用加速器的完整原因链为：图形渲染提供天然大规模数据并行性；游戏市场提供持续销量与研发投入；统一计算设备架构把硬件能力暴露为通用编程模型；科学计算与人工智能又提供了规模更大的规则并行工作负载。四个条件共同形成了今天的图形处理器计算生态。
\end{keybox}
